\documentclass[11pt]{article}
\usepackage[a4paper,width=150mm,top=25mm,bottom=25mm]{geometry}
\usepackage[spanish]{babel}
\usepackage[utf8]{inputenc}
\usepackage{fancyhdr}
\usepackage{xhfill}
\usepackage{nameref}
\pagestyle{fancy}
\usepackage{tikz}
\usetikzlibrary{babel}

\usepackage{blindtext}
\usepackage{graphicx}
\usepackage{tabularx}
\usepackage{listingsutf8}
\usepackage{color}
\usepackage{enumitem}

\definecolor{mygreen}{rgb}{0,0.6,0}
\definecolor{mygray}{rgb}{0.5,0.5,0.5}
\definecolor{mymauve}{rgb}{0.58,0,0.82}
\renewcommand{\lstlistingname}{Listado}
\lstset{
  backgroundcolor=\color{white},   % choose the background color; you must add \usepackage{color} or \usepackage{xcolor}; should come as last argument
  basicstyle=\footnotesize,        % the size of the fonts that are used for the code
  breakatwhitespace=false,         % sets if automatic breaks should only happen at whitespace
  breaklines=true,                 % sets automatic line breaking
  captionpos=b,                    % sets the caption-position to bottom
  commentstyle=\color{mygreen},    % comment style
  deletekeywords={...},            % if you want to delete keywords from the given language
  escapeinside={\%*}{*)},          % if you want to add LaTeX within your code
  firstnumber=1,                   % start line enumeration with line 1
  frame=single,	                   % adds a frame around the code
  inputpath=../servidor/src/sistop/,
  inputencoding=utf8/latin1,
  keepspaces=true,                 % keeps spaces in text, useful for keeping indentation of code (possibly needs columns=flexible)
  keywordstyle=\color{blue},       % keyword style
  language=Java,                 % the language of the code
  morekeywords={*,...},            % if you want to add more keywords to the set
  numbers=none,                    % where to put the line-numbers; possible values are (none, left, right)
  numbersep=5pt,                   % how far the line-numbers are from the code
  numberstyle=\tiny\color{mygray}, % the style that is used for the line-numbers
  rulecolor=\color{black},         % if not set, the frame-color may be changed on line-breaks within not-black text (e.g. comments (green here))
  showspaces=false,                % show spaces everywhere adding particular underscores; it overrides 'showstringspaces'
  showstringspaces=false,          % underline spaces within strings only
  showtabs=false,                  % show tabs within strings adding particular underscores
  stepnumber=2,                    % the step between two line-numbers. If it's 1, each line will be numbered
  stringstyle=\color{mymauve},     % string literal style
  tabsize=2,	                   % sets default tabsize to 2 spaces
  title=\lstname,                   % show the filename of files included with \lstinputlisting; also try caption instead of title
  xleftmargin=15pt,
  xrightmargin=15pt
}
% \usepackage{algorithmicx}
\usepackage{algorithm}
\usepackage{algpseudocode}
% \usepackage{fontawesome}

\fancyhead{}
\fancyhead[R]{Obligatorio}
\fancyhead[L]{Sección \thesection}



\makeatletter
\renewcommand{\ALG@name}{Algoritmo}
\makeatother

\algrenewcommand\algorithmicprocedure{\textbf{procedimiento}}
\algrenewcommand\algorithmicif{\textbf{si}}
\algrenewcommand\algorithmicthen{\textbf{entonces}}
\algrenewcommand\algorithmicelse{\textbf{si no}}
%\algrenewcommand\algorithmicelsif{\textbf{sino si}}
\algrenewcommand\algorithmicend{\textbf{fin}}
\algrenewcommand\algorithmicrequire{\textbf{Pre:}}
\algrenewcommand\algorithmicensure{\textbf{Post:}}
\algrenewcommand\algorithmicforall{\textbf{para cada}}
\algrenewcommand\algorithmicdo{\textbf{hacer}}
\algrenewcommand{\algorithmiccomment}[1]{\hfill$\langle$ #1}

\graphicspath{ {img/} }

\newcommand{\threadFig}[0]{
    \begin{tikzpicture}
      \vspace{5pt}
      \draw[very thick,color=black] plot [smooth,tension=1.5] coordinates{(0.3,0.5) (0.25,0.37) (0.3,0.25) (0.25,0.12)};
    \end{tikzpicture}
}

\newcommand{\BigO}[1]{\ensuremath{\mathrm{O}\bigl(#1\bigr)}}
\newcommand{\type}[1]{\mbox{\textbf{#1}}}
\newcommand{\call}[1]{\mbox{\texttt{#1}}}

\overfullrule=2cm

\begin{document}

\begin{titlepage}
    \begin{center}
        \includegraphics[width=0.7\textwidth]{ucu}

	\vspace*{5cm}

        \Huge
        \textbf{Obligatorio}

        \vspace{0.5cm}
        \LARGE
	Teoría de la Computación

        \vspace{3.5cm}

        Lois, Lucas

        Maresca, Santiago

        Silveira, Marco

        \vfill

        \Large
        Universidad Católica del Uruguay\\
        22 de Noviembre, 2019

    \end{center}
\end{titlepage}



\section{Armado de un Árbol-KDR}
\begin{lstlisting}[language=Python]
	
def makeKDRTree(points, r, dim = 0):
	if not points:
		return None # Empty trees are None
	elif len(points) == 1:
		return tuple(points) # Leaf nodes have one point.

	# Obtenemos las r medianas, para las dimensiones dim, dim+1, ..., dim+(r-1)
	medians = getMedians(points, r, dim)

	partitions = [None] * (2**r + 1) # Tendremos 2**r particiones (hijos)
	partitions[0] = medians
	i = 1
	for criteria in it.product([False, True], repeat = r):  # 2^r iteraciones
		partition = [p for p in points if matches_criteria(criteria, medians, p, dim)] # n iteraciones
		partitions[i] = makeKDRTree(partition, r, (dim+r) % len(points))
		i += 1

	return tuple(partitions)
	\end{lstlisting}
\textbf{ Orden de tiempo de ejecución del algoritmo : }
{Mejor de los casos: points es nulo o es una tupla de largo uno : 
el algoritmo tendrá orden de ejecución : O(1).
Caso promedio : 



}

	\section{Búsqueda de un Árbol-KDR}
	\begin{lstlisting}[language=Python]
	def searchKDRTree(kdrTree, r, point, dim = 0):
		if (kdrTree is None):
			return False
		elif len(kdrTree) == 1:
			return kdrTree[0] == point
		medians = kdrTree[0]
		i = 1
		for criteria in it.product([False, True], repeat = r):  # 2^r iteraciones
			if (matches_criteria(criteria, medians, point, dim)):
				# go to this subtree
				chosen_subtree = kdrTree[i]
				return searchKDRTree(chosen_subtree, r, point, dim+r)
			i += 1
	
		return False
		\end{lstlisting}
	\textbf{ Orden de tiempo de ejecución del algoritmo : }
		{Mejor de los casos: points es nulo o es una tupla de largo uno : 
		el algoritmo tendrá orden de ejecución : O(1).
		Caso promedio : 
		
		
		
		}


% \begin{algorithm}
% \caption{Busca una tupla de largo R en un arbol KDR}
% \label{alg:busqueda-kdr}
% \begin{algorithmic}[1]
% 	\Require El árbol kdr debe estar creado
% \Procedure{searchKDRTree}{kdrTree, r, point, dim = 0}
% \If {kdTree == nulo}
% \State {devolver Falso}
% \Else{}
% \If {kdTree == nulo}
% \State devolver kdrTree[0] == point

% medians =kdrTree[0]
% \ForAll{(criteria,) $\in$ productoCartesiano(false,true,r)}
% 	\State s.put(``\{i\}) \{evento.getNombre()\} \{evento.getCuando()\}'');
% \EndFor
% \EndProcedure
% \end{algorithmic}
% \end{algorithm}

% El algoritmo \ref{alg:crear-kdr} muestra la respuesta del sistema cuando se detalla un id de evento.
% En este caso se representa, para el evento determinado, qué entras están disponibles (mediante id),
%     así como su precio.

\end{document}
